\chapter{系统不变量与权威模型}

\section{核心不变量}

Nous Wave 的各领域和实现必须保持以下关系。

\begin{ReportWideTable}{>{\raggedright\arraybackslash}p{6.1cm}>{\raggedright\arraybackslash}p{8.2cm}}
\captionof{table}{核心设计不变量}\label{tab:invariants}
\toprule
\textbf{不变量} & \textbf{含义}\\
\midrule
主体身份与模型身份分离 & 主体可以更换模型、供应商和运行进程而保持同一逻辑身份。\\
材料、观察事件与认知对象分离 & 原始内容、主体的一次实际接触和长期认知分别拥有独立身份。\\
来源提及与认知指向分离 & 来源中出现某实体，不等于形成的认知必然以该实体为语义对象。\\
证据谱系与独立佐证数量分离 & 多个派生表示如果共享同一来源祖先，只构成同源证据。\\
认知权威、认知运行时与服务投影分离 & 可重建索引、运行时驻留状态和持久认知分别维护。\\
候选提案与权威认知分离 & 模型、离线认知和跨领域推理默认只产生提案，目标领域所有者决定是否提交。\\
检索、呈现、实际使用与长期强化分离 & 被召回或进入上下文不自动增加长期认知支持度。\\
可访问性、相关性、抑制与清除分离 & 回忆难度、当前查询相关程度、普通访问限制和破坏性删除使用不同状态与权限。\\
认知目标与行为承诺分离 & Nous Wave 保存目标和关切；有约束力的承诺、义务和外部执行由 Heptalogos 管理。\\
外部资源与主体认知分离 & 搜索结果和知识库内容只有在主体实际接触并经过认知形成流程后，才可能成为主体认知。\\
\bottomrule
\end{ReportWideTable}

\section{类型化权威}

系统不使用一条覆盖所有语义的全局优先级。不同权威类型分别作用于其负责的事实和操作。

\begin{description}
\item[来源与证据权威] 记录材料内容、观察事件、来源区域、派生过程和可验证外部事实。
\item[认知权威] 由记忆、自我、社会认知和动机等领域维护主体当前接受的认知及其修订。
\item[行为与规范权威] 由 Heptalogos 维护允许、禁止、承诺、义务、行为决策和外部效果。
\item[管理权威] 负责主体生命周期、抑制、恢复、清除、配置和其他显式管理操作。
\end{description}

同一输入可以同时影响多个权威类型，但必须通过各自的领域合同分别处理。例如，一条用户消息首先形成来源事实和观察记录；它可以触发自我或记忆提案，也可以触发 Heptalogos 的行为策略调整。任何一个结果都不能替代其他权威类型的判断。

\section{跨领域提案}

跨领域变化统一通过类型明确的提案传递。提案至少引用目标领域、目标操作、输入对象及其修订、支持证据、生产者身份和形成上下文。目标领域在提交前重新核验来源可用性、当前对象版本、生命周期、访问控制和并发版本。

一个领域不得直接修改另一个领域的权威对象。跨领域原子一致性通过提案引用、提交序号和结果回执建立，而不是通过共享可变对象实现。

\section{管理操作}

自然语言内容不自动获得管理权限。抑制、恢复和清除由显式管理接口和授权上下文触发。普通对话中出现“忘记”“不要再提”等表达时，系统可以记录为观察、偏好或管理建议；破坏性清除仍需独立的管理授权。
